\begin{frame}{왜 ``무작위 걷기''가 ``문장''을 대신하는 자연스러운 선택인가}
  \begin{itemize}
    \item word2vec에서 ``문장''이 가진 본질: \textbf{어떤 것들이 어떤
      순서로 함께 나타나는지}를 담는 것.
    \item 그래프 노드의 ``함께 나타남'' = ``엣지로 연결되어,
      서로를 따라다니면 만나는 것''.
    \item 어떤 노드에서 출발해 이웃을 무작위로 골라 반복 이동하는
      \textbf{무작위 걷기}가 바로 이 ``함께 나타남의 순서''를 만들어내는
      \textbf{최소 장치}.
    \item 걷기 안에서 \textbf{인접한 노드는 서로 엣지로 연결되어} 있고,
      어떤 노드 주위에 자주 등장하는 노드 = 그 노드의 구조적 ``이웃''.
  \end{itemize}
\end{frame}
